\section{Marco Teórico}

\subsection{Computación de Alto Rendimiento (HPC)}

La Computación de Alto Rendimiento (HPC, por sus siglas en inglés) consiste en la
agregación de potencia de cálculo a través de arquitecturas paralelas para resolver
problemas científicos y tecnológicos de extrema complejidad \parencite{hager2010hpc}. Esta
disciplina permite ejecutar algoritmos y modelados matemáticos que resultarían
inabarcables para un sistema tradicional debido a restricciones de tiempo y recursos de
hardware. Al dividir cargas masivas de trabajo en instrucciones más pequeñas que se
procesan simultáneamente, HPC maximiza la velocidad de procesamiento de datos y la
obtención de resultados en la ingeniería \parencite{patterson2017computer}.

Históricamente, el aumento de la frecuencia de reloj en los microprocesadores permitía
acelerar la ejecución de algoritmos secuenciales de forma directa. Sin embargo, tras
alcanzar los límites físicos y térmicos del silicio, la industria se desplazó hacia
arquitecturas multinúcleo que exigen un paradigma de programación paralela
\parencite{patterson2017computer}. Bajo este enfoque, el diseño de software eficiente requiere
distribuir la carga de trabajo de manera concurrente, haciendo del análisis de algoritmos
iterativos numéricos, como el método de Jacobi, un caso de estudio en HPC.

\subsection{La Ecuación de Poisson en 1D}

La ecuación de Poisson en una dimensión es una ecuación diferencial de segundo orden de
la forma:
\begin{equation}
    -\frac{d^2u}{dx^2} = f(x)
    \label{eq:poisson_1d}
\end{equation}
definida sobre un dominio cerrado $[a, b]$, donde $f(x) \in C[a, b]$ es una función de
forzamiento dada \parencite{burkardt2011jacobi}. Para garantizar la unicidad de la solución,
esta ecuación se complementa con condiciones de frontera de Dirichlet en los extremos del
dominio:
\begin{equation}
    u(a) = u_a, \qquad u(b) = u_b
    \label{eq:dirichlet}
\end{equation}
Un ejemplo concreto y representativo se define sobre el intervalo $[0, 1]$ con condiciones
$u(0) = u(1) = 0$, función de forzamiento $f(x) = -x(x+3)e^x$ y solución analítica
exacta $u(x) = x(x-1)e^x$.

La ecuación de Poisson es la representación estacionaria de la ecuación de calor y
aparece de forma ubícua en simulación de fluidos, electrostática, mecánica de sólidos y
transferencia de calor \parencite{elman2005finite}. Su relevancia en HPC radica en que su
resolución numérica, especialmente para dominios de alta dimensión o mallas finas, genera
sistemas lineales de gran escala cuyo costo computacional motiva el uso de métodos
iterativos y estrategias de paralelismo \parencite{burkardt2011jacobi}.

\subsection{Discretización del Operador de Poisson}

La solución numérica de la ecuación de Poisson requiere reemplazar el operador
diferencial continuo por una aproximación discreta sobre una malla de $N$ nodos $x_j$
distribuidos uniformemente en $[a, b]$ con espaciado $h = (b-a)/(N-1)$
\parencite{burkardt2011jacobi}. Aplicando diferencias finitas centradas, la segunda derivada en
el nodo $x_j$ se aproxima mediante el cociente de diferencias:
\begin{equation}
    -\frac{d^2u}{dx^2}(x_j) \approx \frac{-u_{j-1} + 2u_j - u_{j+1}}{h^2}
    \label{eq:fd_stencil}
\end{equation}
Esta expresión es una aproximación de segundo orden al operador de Poisson y constituye
la ecuación discreta en cada nodo interior de la malla \parencite{burkardt2011jacobi}. Para
garantizar que los errores de discretización disminuyan apropiadamente al refinar la
malla, se selecciona una secuencia anidada de grillas con $n_k = 2^k + 1$ nodos, de
forma que el espaciado $h_k = (b-a)/2^k$ converge a cero conforme $k \to \infty$.

\subsection{El Sistema Lineal Discretizado}

Al aplicar las ecuaciones de diferencias finitas en todos los nodos interiores, junto con
las condiciones de frontera de Dirichlet, se obtiene un sistema lineal
$A \cdot \mathbf{u} = \mathbf{f}$ de dimensión $N \times N$ \parencite{burkardt2011jacobi}.
La matriz del sistema $h^2 A$ es tridiagonal, simétrica y definida positiva, con la
estructura:
\begin{equation}
    h^2 A =
    \begin{pmatrix}
        1  & 0  & 0  & 0  & \cdots & 0 \\
        -1 & 2  & -1 & 0  & \cdots & 0 \\
        0  & -1 & 2  & -1 & \cdots & 0 \\
        \vdots & & \ddots & \ddots & \ddots & \vdots \\
        0  & 0  & \cdots & -1 & 2  & -1 \\
        0  & 0  & \cdots & 0  & 0  & 1
    \end{pmatrix}
    \label{eq:matriz_tridiagonal}
\end{equation}
Las propiedades de simetría y definición positiva son condiciones suficientes para
garantizar la convergencia de los métodos iterativos clásicos, en particular la iteración
de Jacobi \parencite{burkardt2011jacobi}. Este sistema presenta una estructura dispersa
(\textit{sparse}) altamente favorable desde el punto de vista computacional, ya que cada
fila contiene a lo sumo tres elementos no nulos, lo que reduce significativamente el costo
de memoria y operaciones por iteración \parencite{barrett1994templates}.

\subsection{El Método Iterativo de Jacobi}

Dado un sistema lineal $A \cdot \mathbf{x} = \mathbf{f}$ y una aproximación inicial
$\mathbf{x}^{(0)}$, la iteración de Jacobi genera una secuencia de aproximaciones
$\mathbf{x}^{(0)}, \mathbf{x}^{(1)}, \mathbf{x}^{(2)}, \ldots$ mediante la actualización
simultánea de todos los componentes de la solución \parencite{burkardt2011jacobi}. La
actualización del $i$-ésimo componente en la iteración $(k+1)$ se define como:
\begin{equation}
    x_i^{(k+1)} = \frac{f_i - \displaystyle\sum_{j \neq i} A_{ij} \cdot x_j^{(k)}}{A_{ii}}
    \label{eq:jacobi_scalar}
\end{equation}
Bajo la descomposición matricial $A = L + D + U$, donde $L$, $D$ y $U$ son las partes
triangular inferior, diagonal y triangular superior respectivamente, la iteración se
expresa de forma vectorizada como:
\begin{equation}
    \mathbf{x}^{(k+1)} = D^{-1}\bigl(\mathbf{f} - (L + U)\,\mathbf{x}^{(k)}\bigr)
    \label{eq:jacobi_vectorized}
\end{equation}
Esta formulación vectorizada es computacionalmente eficiente, pues la actualización de
cada componente en una iteración es completamente independiente de los demás, propiedad
que resulta fundamental para explotar el paralelismo de datos mediante hilos y procesos
\parencite{burkardt2011jacobi}.

\subsection{Convergencia del Método de Jacobi}

El criterio de parada de la iteración de Jacobi se basa en el monitoreo del residual
$\mathbf{r}^{(k)} = A \cdot \mathbf{x}^{(k)} - \mathbf{f}$ \parencite{burkardt2011jacobi}.
La métrica estándar empleada es la norma RMS del residual:
\begin{equation}
    \text{Residual RMS} = \frac{\|\mathbf{r}^{(k)}\|}{\sqrt{N}} \leq \varepsilon
    \label{eq:residual_rms}
\end{equation}
que permite comparar la convergencia entre sistemas de diferentes tamaños de malla con
una tolerancia fija $\varepsilon$ \parencite{burkardt2011jacobi}. Un aspecto crítico del método
de Jacobi aplicado al problema de Poisson es su convergencia lenta: el número de
iteraciones $m_k$ requeridas crece aproximadamente como $m_k \sim \mathcal{O}(n_k^2)$,
implicando que el costo computacional total escala como $\mathcal{O}(n_k^3)$. 
Esta degradación severa al refinar la malla motiva directamente la necesidad de paralelismo: al 
distribuir las iteraciones entre múltiples hilos o procesos, es posible reducir el tiempo de 
ejecución real incluso cuando el número total de operaciones aritméticas permanece constante \parencite{briggs2000multigrid}.

\subsection{Programación Secuencial y Concurrente}

La programación secuencial es el paradigma clásico donde las instrucciones de un
algoritmo se ejecutan de forma estrictamente lineal sobre un único núcleo de
procesamiento \parencite{benari2006principles}. En este modelo, el límite de rendimiento está
directamente dictado por la frecuencia de reloj del hardware y el número de iteraciones
del algoritmo. Como respuesta a estas limitaciones, la programación concurrente estructura
el software en múltiples hilos o tareas independientes que hacen progreso de forma
superpuesta en el tiempo, reduciendo drásticamente el tiempo de ejecución en problemas de
cómputo intensivo.

Para el algoritmo de Jacobi, la paralelización es conceptualmente directa: dado que la
actualización de cada nodo $x_i^{(k+1)}$ depende únicamente de los valores de la
iteración anterior $\mathbf{x}^{(k)}$, el cálculo de todos los nodos interiores en una
misma iteración es completamente independiente entre sí \parencite{burkardt2011jacobi}. Esta
independencia de datos por iteración elimina las condiciones de carrera durante la fase de
actualización, permitiendo distribuir los nodos entre hilos o procesos sin necesidad de
sincronización intra-iteración, aunque sí se requiere una barrera de sincronización entre
iteraciones consecutivas \parencite{silberschatz2018os}.

\subsection{Hilos (Threads) y la Librería Pthreads}

En el ámbito de la computación paralela, el modelo de memoria compartida establece que
múltiples hilos de ejecución creados por la librería POSIX Threads (Pthreads) operan
dentro del mismo espacio de direccionamiento lógico del proceso padre
\parencite{silberschatz2018os}. A nivel de arquitectura, esto significa que todos los núcleos
del procesador acceden de forma directa y simultánea a las mismas estructuras de datos
residentes en la memoria principal, sin necesidad de copiar o transmitir datos entre
unidades de procesamiento. Para el algoritmo de Jacobi, este modelo es altamente
eficiente, pues el vector de solución $\mathbf{u}$ y el vector de forzamiento $\mathbf{f}$
se comparten íntegramente entre todos los hilos, y cada uno opera únicamente sobre un
subconjunto disjunto de nodos, eliminando la contención de escritura
\parencite{burkardt2011jacobi}.

La creación de un hilo mediante \texttt{pthread\_create} y su sincronización mediante
\texttt{pthread\_join} introduce una sobrecarga fija por hilo que resulta despreciable
cuando el número de nodos por hilo es suficientemente grande \parencite{silberschatz2018os}.
No obstante, entre iteraciones de Jacobi es imprescindible garantizar que todos los hilos
hayan completado la actualización de sus nodos antes de iniciar la siguiente iteración;
esto se logra mediante una barrera de sincronización (\texttt{pthread\_barrier\_wait}),
cuyo costo adicional es un factor determinante en el análisis de escalabilidad
\parencite{hager2010hpc}.

\subsection{Procesos y la Llamada al Sistema Fork}

La paralelización mediante la llamada al sistema \texttt{fork()} crea procesos hijos
independientes que no comparten espacio de memoria con el proceso padre, operando bajo
el modelo de memoria distribuida dentro de un mismo nodo físico \parencite{silberschatz2018os}.
Cada proceso hijo posee su propia copia privada del espacio de direcciones, obtenida
mediante el mecanismo de \textit{copy-on-write} del sistema operativo, lo que elimina las
condiciones de carrera pero introduce una sobrecarga significativa en la creación, gestión
y comunicación entre procesos. Para el algoritmo de Jacobi,
cada proceso hijo calcula la actualización de un subconjunto de nodos, pero los resultados
parciales deben comunicarse al proceso padre mediante mecanismos de IPC
(\textit{Inter-Process Communication}) como tuberías (\textit{pipes}) o memoria
compartida explícita (\texttt{shmget}/\texttt{mmap}) antes de iniciar la siguiente
iteración.

La principal desventaja del modelo basado en \texttt{fork} para algoritmos iterativos
como Jacobi reside en el costo acumulado de esta comunicación inter-iteración
\parencite{silberschatz2018os}. Dado que el método de Jacobi puede requerir miles de
iteraciones para converger, especialmente en mallas finas, la sobrecarga de
sincronizar resultados entre procesos en cada iteración puede superar ampliamente la
ganancia obtenida por la distribución del cómputo, comportamiento análogo al evidenciado
previamente en la paralelización por procesos de la multiplicación de matrices
\parencite{hager2010hpc}.

\subsection{Jerarquía de Caché y Localidad de Memoria}

El rendimiento de una aplicación HPC no depende exclusivamente de la capacidad matemática
del procesador, sino también del acceso al subsistema de memoria \parencite{drepper2007memory}.
Los procesadores modernos integran una jerarquía de memoria ultrarrápida (L1, L2 y L3)
cuya unidad mínima de transferencia es la \textit{cache line} de 64 bytes. Para el algoritmo de Jacobi en 1D, el patrón de acceso al
vector de solución $\mathbf{u}$ es estrictamente secuencial (nodo $j-1$, $j$, $j+1$),
lo que favorece de forma natural la localidad espacial de memoria y maximiza el
aprovechamiento de la caché, a diferencia de los accesos por columnas del caso matricial
\parencite{patterson2017computer}.

Sin embargo, en la implementación paralela, la división del vector entre múltiples hilos
introduce el fenómeno de \textit{false sharing}: si dos hilos adyacentes operan sobre
nodos que residen en la misma línea de caché, las escrituras de un hilo invalidan la
línea en el caché del otro, generando una contención invisible a nivel de código que
penaliza severamente el rendimiento. Garantizar que cada hilo
opere sobre bloques de nodos suficientemente grandes, y alineados con el tamaño de la
línea de caché, es una condición necesaria para minimizar este efecto en la
implementación paralela del método de Jacobi \parencite{drepper2007memory}.

\subsection{Métricas de Rendimiento en Computación Paralela}

Para cuantificar el éxito de una implementación paralela, la literatura científica emplea
métricas matemáticas estandarizadas \parencite{hager2010hpc}. La métrica fundamental es el
\textit{Speedup} o Aceleración $S_p$, definido como:
\begin{equation}
    S_p = \frac{T_1}{T_p}
    \label{eq:speedup}
\end{equation}
donde $T_1$ es el tiempo de ejecución del algoritmo secuencial óptimo y $T_p$ es el
tiempo paralelo con $p$ unidades de procesamiento \parencite{hager2010hpc}. Derivada de esta
métrica, la Eficiencia $E_p$ evalúa el grado de aprovechamiento de los recursos físicos:
\begin{equation}
    E_p = \frac{S_p}{p} = \frac{T_1}{p \cdot T_p}
    \label{eq:eficiencia}
\end{equation}

El desempeño asintótico de estos sistemas está regido por la \textbf{Ley de Amdahl}, que
establece que el Speedup máximo está limitado por la fracción no paralelizable $(1-f)$
del algoritmo:
\begin{equation}
    S_{\text{Amdahl}} \leq \frac{1}{(1-f) + \dfrac{f}{p}}
    \label{eq:amdahl}
\end{equation}
Para el algoritmo de Jacobi, la fracción secuencial incluye la evaluación del residual
global de convergencia y la barrera de sincronización entre iteraciones, que no pueden
paralelizarse y constituyen el cuello de botella principal según este modelo
\parencite{patterson2017computer}. En contraste, la \textbf{Ley de Gustafson-Barsis} describe
la escalabilidad débil y argumenta que, al incrementar el tamaño de la malla $N$, la
fracción paralelizable del trabajo crece proporcionalmente, haciendo que la porción
secuencial pierda peso relativo:
\begin{equation}
    S_{\text{Gustafson}} = (1-f) + p \cdot f
    \label{eq:gustafson}
\end{equation}
Bajo la visión de Gustafson, a medida que la malla se refina, la cantidad de operaciones
paralelas aumenta drásticamente, mitigando el cuello de botella pronosticado por Amdahl y
permitiendo un Speedup escalable con el número de unidades de procesamiento
\parencite{gustafson1988reevaluating}.
